%% -*- coding:utf-8 -*-

% general stuff for grammar theory is in local packages.tex
% everything that is specific to Chinese is in here.

%% below are added by Lulu for the Chinese version of the book
\usepackage{zhnumber}
% to type­set Chinese representations of numbers.  

%\usepackage[UTF8, heading = true]{ctex}

\usepackage{titlesec}  
\AtBeginDocument{

% for Chinese style of titles of chapters and sections.
\renewcommand{\partname}{第\thepart{}部分}  
\def\thepart{\zhdig{part}}
%\pagenumbering {zhnum}
\renewcommand{\chaptername}{第\zhnum{chapter}章}  
\renewcommand{\appendixname}{附录}
\appto\appendix{
\renewcommand{\chaptername}{\appendixname\thechapter} 
}
%\renewcommand{\chapterautorefname}{第\zhnum{chapter}章}%

% for Chinese style of pages. 15.4.2019
%   \def\reftextfaceafter{在\reftextvario{下}{后} 一页}%
%    \def\reftextfacebefore{在\reftextvario{上}{前} 一页}%
%    \def\reftextafter{在前一页}%
%    \def\reftextbefore{在后一页}%
%    \def\reftextcurrent{在这一页}%
%    \def\reftextfaraway#1{在第~\pageref{#1}页}%
%    \def\reftextpagerange#1#2{在第~\pageref{#1}--\pageref{#2}页}%
%    \def\reftextlabelrange#1#2{从第\ref{#1}页到第~\ref{#2}页}%

% The following would also change the chapter and section number in the heading. We do not want
% this. 27.10.2016
%\def\thechapter{\zhnum{chapter}}
%\def\thesection{\zhnum{section}}

%\lsBookLanguageChinese
%\renewcommand{\lsBookLanguage}{chinese}

%re\newcommand{\sectionname}{节}  
\renewcommand{\figurename}{图}  
\renewcommand{\tablename}{表}  
%\renewcommand{\bibname}{参考文献}  
\renewcommand{\contentsname}{目~~~~录}  
\renewcommand{\listfigurename}{图~目~录}  
\renewcommand{\listtablename}{表~目~录}  
\renewcommand{\indexname}{索~引}  
\renewcommand{\abstractname}{\Large{摘~要}}  
\newcommand{\keywords}[1]{\\ \\ \textbf{关~键~词}：#1}  
\titleformat{\chapter}[block]{\center\Large\bf}{\chaptername}{20pt}{}  
\titleformat{\appendix}[block]{\center\Large\bf}{\appendixname}{20pt}{}  
\titleformat{\section}[block]{\large\bf}{\thesection}{10pt}{}  
\titleformat{\tableofcontents}[block]{\center\Large\bf}{\contensname}{20pt}{}  
\titleformat{\part}[block]{\center\Large\bf}{\partname}{20pt}{}  
\titleformat{\index}[block]{\center\Large\bf}{\indexname}{20pt}{}  
}
% for Chinese style of table of contents
%\titlecontents{\chapter}

\usepackage{indentfirst}
% In Chinese, each paragraph should be indented.
\setlength{\parindent}{2em}
% the indented space should be two characters.

% quotes should be in simkai（楷体）
\newCJKfontfamily[\fontpath/halfish simkai]\simkai{KaiTi}[Mapping=tex-text,Ligatures=Common,Scale=MatchUppercase]

% indentation of the whole quotation is 2em like the paragraph indentation
% St. Mü. 22.04.2018
\makeatletter
\renewenvironment{quotation}
			   {\list{}{\leftmargin2em\listparindent 2em%
						%\itemindent    \listparindent
						%\rightmargin   \leftmargin
						\parsep        \z@ \@plus\p@}%
				\item\relax}
			   {\endlist}
\makeatother
\AtBeginEnvironment{quotation}{\simkai}

% this influences enumerate and itemize
\RequirePackage{enumitem}
\setlist{leftmargin=3em}


%\usepackage{newunicodechar}
%\newunicodechar{(}{（}



% use Chinese parenthesis
% https://tex.stackexchange.com/a/428832/18561
\makeatletter
\renewcommand*{\bibleftparen}{\unspace\blx@postpunct（}
%\renewcommand*{\bibleftparen}{\blx@postpunct（}
\renewcommand*{\bibrightparen}{\blx@postpunct）\unspace\midsentence}
%\renewcommand*{\bibrightparen}{\blx@postpunct）\midsentence}




\AtBeginDocument{
\newindex{sbc}{scx}{scd}{中文术语索引}
}

\newcommand{\footnoteindex@sbc}[2]{\index[sbc]{#2|infn{#1}}}
\newcommand{\isc}[1]{%
  \if@noftnote%
    \index[sbc]{#1}%
    \else%
    \edef\tempnumber{\thefootnote}%
    \expandafter\footnoteindex@sbc\expandafter{\tempnumber}{#1}%
    %\indexftn{#1}{\value{footnotemark}}%
  \fi%
}


\renewcommand{\lsImpressum}{
\raggedright

\lsImpressumCitationText

\vfill

本书可在以下网址获取： \\
%This title can be downloaded at:\\
\url{\lsURL}

© \the\year, \iflsCollection the authors\else\@author\fi

\newcommand{\ccby}{CC-BY} 
\ifx\lsCopyright\ccby 
本书的发表采用创作共用许可证4.0（CC BY 4.0）授权：
%Published under the Creative Commons Attribution 4.0 Licence (CC BY 4.0):
http://creativecommons.org/licenses/by/4.0/ 
\else
本书的发表采用创作共用许可证—禁止演绎4.0（CC BY-ND 4.0）授权：
%Published under the Creative Commons Attribution-NoDerivatives 4.0 Licence (CC BY-ND 4.0):
http://creativecommons.org/licenses/by-nd/4.0/ 
\fi

\noindent 
\begin{tabular}{@{}llrl@{; }r@{~}l}
ISBN: & \multicolumn{2}{r}{完整电子版：} & \lsISBNdigital \\
      &	精装：& vol1:        & \lsISBNhardcoverOne & vol. 2 & \lsISBNhardcoverTwo\\
      &	简装：& vol1:        & \lsISBNsoftcoverOne & vol. 2 & \lsISBNsoftcoverTwo\\
      &	美版简装：& vol1:        & \lsISBNsoftcoverusOne & vol. 2 & \lsISBNsoftcoverusTwo \\
%ISBN: & \multicolumn{2}{r}{Digital, complete work:} & \lsISBNdigital \\
%      &	Hardcover: & vol1:        & \lsISBNhardcoverOne & vol. 2 & \lsISBNhardcoverTwo\\
%      &	Softcover: & vol1:        & \lsISBNsoftcoverOne & vol. 2 & \lsISBNsoftcoverTwo\\
 %     &	Softcover US: & vol1:        & \lsISBNsoftcoverusOne & vol. 2 & \lsISBNsoftcoverusTwo \\
\end{tabular}

\IfStrEq{\lsISSN}{??}		% \IfStrEq from xstring
	{}
	{ISSN: \lsISSN}

\IfStrEq{\lsBookDOI}{??}		% \IfStrEq from xstring
	{}
	{\doi{\lsBookDOI}}


\bigskip

封面设计： 
%Cover and concept of design:
Ulrike Harbort \\
\if\@translator\empty\else
翻译：
%Translators:
\@translator \\
\fi
\if\@typesetter\empty\else
排版：
%Typesetting:
\@typesetter \\
\fi
\if\@illustrator\empty\else
插图：
%Illustration:
\@illustrator \\
\fi
\if\@proofreader\empty\else
校对：
%Proofreading:
\@proofreader \\
\fi
\if\@openreviewer\empty\else
公开评审：
%Open reviewing:
\@openreviewer \\
\fi
字体：Linux Libertine, Arimo, DejaVu Sans Mono\lsAdditionalFontsImprint\\
%Fonts: Linux Libertine, Arimo, DejaVu Sans Mono\lsAdditionalFontsImprint\\
排版软件：\XeLaTeX
%Typesetting software: \XeLaTeX


\bigskip

%语言科学出版社\\
Language Science Press\\
Unter den Linden 6\\
10099 Berlin, Germany\\
\href{http://langsci-press.org}{langsci-press.org}

\vfill

本书由柏林自由大学存储和编目 \\[3ex]
%Storage and cataloguing done by FU Berlin \\[3ex]  

\IfFileExists{\logopath/FULogo_sw_CMYK.pdf}{\includegraphics[width=5cm]{\logopath/FULogo_sw_CMYK.pdf}}{FU-Logo} \\[3ex]

\vfill

\noindent
本书的出版机构语言科学出版社不负责保证外部或第三方网址保持不变或者准确性，也不能保证这些网站内容是（或者保持）准确和合理的。
%\lsp has no responsibility for the persistence or accuracy of URLs for
%external or third-party Internet websites referred to in this
%publication, and does not guarantee that any content on such websites
%is, or will remain, accurate or appropriate. 
%Information regarding prices, travel timetables and other factual information given in this work are correct at the time of first publication but \lsp does not guarantee the accuracy of such information thereafter. 
}


\makeatother



\newcommand{\mytrans}[1]{\trans\quotetrans{#1}}

% use Chinese single quotes and adjust the kerning
\newcommand{\quotetrans}[1]{‘#1’}
\xeCJKsetkern{。}{’}{0.2em}

% The following does not work since we have trnaslations ending in question marks or without punctuation
%\newcommand{\quotetrans}[1]{`#1\hspace{-0.2em}'}



% use a less heavy variant of Libertine: Display font
  %% \setmainfont[
  %%         Ligatures={TeX,Common},
  %%         Path=\fontpath,
  %%         PunctuationSpace=0,
  %%         Numbers={Proportional},
  %%         BoldFont = LinLibertine_RZ.otf ,
  %%         ItalicFont = LinLibertine_RI.otf ,
  %%         BoldItalicFont = LinLibertine_RZI.otf,
  %%         BoldSlantedFont = LinLibertine_RZ.otf,
  %%         SlantedFont    = LinLibertine_R.otf,
  %%         SlantedFeatures = {FakeSlant=0.25},
  %%         BoldSlantedFeatures = {FakeSlant=0.25},
  %%         SmallCapsFeatures = {FakeSlant=0,Numbers=OldStyle},
  %%         ]{LinLibertine_DR.otf} 



%\setCJKmainfont{LinLibertine_R.otf}


\makeatletter
\newlength\fake@f
\newlength\fake@c
\def\fakesc#1{%
  \begingroup%
  \xdef\fake@name{\csname\curr@fontshape/\f@size\endcsname}%
  \fontsize{\fontdimen8\fake@name}{\baselineskip}\selectfont%
  \uppercase{#1}%
  \endgroup%
}
\makeatother



%\setmainfont{SourceHanSerifSC-ExtraLight}\setCJKmainfont{SourceHanSerifSC-ExtraLight}  %\let\textsc\fakesc
%\setmainfont{SourceHanSerifSC-Light}\setCJKmainfont{SourceHanSerifSC-Light}  %\let\textsc\fakesc
%\setmainfont{SourceHanSerifSC-Medium}\setCJKmainfont{SourceHanSerifSC-Medium}  %\let\textsc\fakesc


%\setmainfont{cmunrm.otf}

%\setmainfont{SimSun}\setCJKmainfont{SimSun}

%\setCJKmainfont{SourceHanSerifSC-ExtraLight} zu dünn

%\setCJKmainfont{SourceHanSerifSC-Light} zu dünn

%\setCJKmainfont{SourceHanSerifSC-Medium}  nicht schön ???


%\setCJKmainfont[Scale=MatchUppercase]{SourceHanSerifSC-Regular}


%\setmainfont{SourceSerifPro-Regular}\setCJKmainfont[Scale=MatchUppercase]{SourceHanSerifSC-Regular}

% Download fonts at: https://github.com/adobe-fonts/source-han-serif
% https://github.com/adobe-fonts/source-serif-pro

\setmainfont[BoldFont={SourceSerifPro-Bold}, 
 ItalicFont={SourceSerifPro-LightIt},
 BoldItalicFont={SourceSerifPro-BoldIt}]{SourceSerifPro-Light}\setCJKmainfont[Scale=MatchUppercase]{SourceHanSerifSC-Regular}

\newcommand{\passive}{被}
\newcommand{\passivepst}{\passive.\textsc{pst}}
\newcommand{\passiveprs}{\passive.\textsc{prs}}
\newcommand{\passiveimp}{\passive.\textsc{imp}}

